\section{Leading terms and transfer of jets}\label{sec:transfer}

We first turn spaces of sections into finite sets of exponents.  The
essential point is that jet separation by the resulting monomials implies
jet separation by the original sections.  We will also make two elementary
changes to the exponent sets while retaining this implication.
These operations follow the support method of
\cite[Section~2]{OpenAI2026Higher}; the formal transfer below makes the
method available in arbitrary characteristic.

Throughout, $k$ is algebraically closed.  The arguments of this section
work in every characteristic.  For a smooth $d$-dimensional variety $Y$,
a line bundle $A$, and $y\in Y(k)$, the jets through order $m\geq0$ are
\[
 J_y^m(A)=A_y/\mathfrak m_y^{m+1}A_y.
\]
A finite-dimensional space $W\subset H^0(Y,A)$ \emph{separates these
jets at} $y_1,\ldots,y_r$ if the evaluation map
$W\to\bigoplus_i J_{y_i}^m(A)$ is surjective.  The locus of such
ordered distinct tuples is open in $Y^r$: the jet spaces are the fibers
of the locally free sheaf of principal parts of $A$, of rank
$\binom{d+m}{d}$, and surjectivity is a rank condition.  The description
of principal parts by the $m$th infinitesimal neighborhood of the diagonal
also shows compatibility with field extension.  Local coordinates for a
smooth variety identify its fibers with truncated polynomial rings.
For the diagonal description and base change, see
\cite[Tags~0H90 and~0H8Z]{StacksProject}.  Local freeness follows as well
from the filtration whose graded pieces are
$A\otimes\operatorname{Sym}^j\Omega_Y^1$, $0\leq j\leq m$.

Put $\T^d=(\mathbb G_{m,k})^d$.  For a finite set $S\subset\Z^d$, write
\[
 M(S)=\Span_k\{x^\alpha:\alpha\in S\}
 \subset k[x_1^{\pm1},\ldots,x_d^{\pm1}].
\]
Here and below, auxiliary existence statements about points on a torus
refer to $k$-points.

\begin{definition}\label{def:graded}
A family of finite sets $(S_N)_{N\geq1}$ in $\Z^d$ is \emph{graded} if
\[
 S_N+S_{N'}\subset S_{N+N'}\qquad(N,N'\geq1).
\]
A transformation $(S_N)\mapsto(T_N)$ has \emph{backward jet transfer}
if, for every $N,r\geq1$ and $m\geq0$, separation of jets through order
$m$ by $M(T_N)$ at some $r$ distinct torus points implies the same
assertion for $M(S_N)$, possibly at a different tuple.
\end{definition}

\subsection{Least exponents}

A monomial order on $\Z_{\geq0}^h$ means an addition-compatible total
well-order.  For a nonzero formal series $f\in k[[z_1,\ldots,z_h]]$,
let $\nu(f)$ be the least exponent occurring in $f$, for a fixed such
order.  In particular,
\begin{equation}\label{eq:least-product}
 \nu(fg)=\nu(f)+\nu(g).
\end{equation}
Indeed, the product of the two least terms is the unique least term in
the product.

\begin{lemma}\label{lem:least-terms}
For a finite-dimensional subspace $W\subset k[[z_1,\ldots,z_h]]$, the
set $\{\nu(f):0\neq f\in W\}$ has cardinality $\dim_k W$.
More generally, suppose finite-dimensional spaces $W_{N,\gamma}$, with
only finitely many nonzero spaces in each degree, are indexed by
$N\geq1$ and $\gamma\in\Z^{d-h}$ and satisfy
\[
 W_{N,\gamma}W_{N',\gamma'}
 \subset W_{N+N',\gamma+\gamma'}.
\]
Then recording the pairs $(\nu(f),\gamma)$ for $0\neq f\in W_{N,\gamma}$
gives a graded family, with cardinality $\sum_\gamma\dim_k W_{N,\gamma}$
in degree $N$.
\end{lemma}

\begin{proof}
Choose the smallest exponent occurring as $\nu(f)$ for a nonzero member
of $W$, and choose a member whose coefficient there is $1$.  The kernel
of that coefficient functional has codimension one.  Repeating in this
kernel gives a basis with distinct least exponents.  Every nonzero linear
combination has as its least exponent the smallest least exponent among
the basis vectors with nonzero coefficients.  This proves the first
assertion; the second follows label by label and from
\eqref{eq:least-product}.
\end{proof}

For the transfer of jets we need to realize finitely many chosen least
terms by a single positive weight vector.  The weights may depend on the
degree; the monomial order defining the graded family will remain fixed.

\begin{lemma}\label{lem:exposing-weights}
Let $f_1,\ldots,f_q$ be nonzero formal series, and let $e_j=\nu(f_j)$.
There is a vector $\lambda\in\Z_{>0}^h$ for which $e_j$ is the unique
exponent of least $\lambda$-weight in $f_j$, simultaneously for all $j$,
in either of the following cases:
\begin{enumerate}[label=\textup{(\roman*)}]
\item the order is lexicographic;
\item $h=2$ and, for a fixed real $t>0$, the order compares
$(\alpha+t\beta,\beta)$ lexicographically.
\end{enumerate}
\end{lemma}

\begin{proof}
In the lexicographic case set $\lambda_h=1$, and successively choose
\[
 \lambda_i>\max_j\sum_{a>i}\lambda_a(e_j)_a
 \qquad(i=h-1,\ldots,1).
\]
If an exponent is lexicographically greater than $e_j$, its gain at the
first differing coordinate outweighs every possible loss in the later
coordinates.

For (ii), the indicated order is a monomial order because every bounded
weight range contains only finitely many nonnegative integral exponents.
Choose a real bound larger than the weights of all the $e_j$.  Among the
finitely many exponents below that bound, replacing $t$ by a sufficiently
close rational $t'>t$ preserves all strict comparisons with the $e_j$
and resolves equal-weight comparisons by increasing $\beta$.  Exponents
above the bound cannot interfere: their weights only increase, while the
new weights of the $e_j$ stay below the bound.  Clearing the denominator
of $(1,t')$ gives the required positive integral weights.
\end{proof}

\subsection{Formal specialization of jet matrices}

The next lemma supplies the passage from formal leading terms to actual
points.  Ito's weighted degeneration argument gives such a jet-rank
comparison over $\mathbb C$ \cite[Proposition~5.7]{Ito2013}.  Here
substitution in truncated rings supplies a formal version in arbitrary
characteristic, without differentiation or factorial denominators.

\begin{lemma}[Formal jet transfer]\label{lem:formal-transfer}
Let $1\leq h\leq d$, let $Y$ be a smooth integral $h$-dimensional
$k$-variety, and choose $y\in Y(k)$ with formal coordinates
$z_1,\ldots,z_h$.  Let $f_1,\ldots,f_q$ be regular functions on a common
neighborhood of $y$, and let $\gamma_j\in\Z^{d-h}$.
Suppose the expansion of $f_j$ has a term $c_jz^{e_j}$, with
$c_j\neq0$, which is uniquely lowest for one common weight vector
$\lambda\in\Z_{>0}^h$.

If the span of the monomials $Z^{e_j}u^{\gamma_j}$ separates jets
through order $m$ at $r$ distinct points of $\T^h\times\T^{d-h}$,
then the span of $f_ju^{\gamma_j}$ separates the same jets at $r$
distinct points in the product of that neighborhood with $\T^{d-h}$.
The points may be chosen in any prescribed dense open subset of this
product.  The statement also holds for sections written in a local
frame of a line bundle.
\end{lemma}

\begin{proof}
Let $(Q_i,U_i)$, $1\leq i\leq r$, be a tuple at which the monomials
separate the stated jets.  Shrink the neighborhood of $y$ to an affine
one, denoted $Y^\circ$, on which all the functions are defined.  Put
\[
 R=k[[s]],\qquad F=k((s)),\qquad
 B=R[\delta_1,\ldots,\delta_d]/(\delta_1,\ldots,\delta_d)^{m+1},
\]
and write $\delta=(\delta',\delta'')$ according to the two factors.
For each $i$, substitute
\begin{equation}\label{eq:formal-substitution}
 z_a=s^{\lambda_a}(Q_{i,a}+\delta'_a),\qquad
 u_b=U_{i,b}+\delta''_b.
\end{equation}
The first substitution is defined on the completed local ring, since
$B$ is $s$-adically complete and every $\lambda_a$ is positive.  The
second is defined on Laurent polynomials, since each $U_{i,b}$ is
nonzero.  In particular, these substitutions define ring maps from the
coordinate ring of $Y^\circ\times\T^{d-h}$ to $B$.

Use the classes of $\delta^\alpha$, $|\alpha|\leq m$, as an $R$-basis
of $B$.  The images of $f_ju^{\gamma_j}$ under all $r$ substitutions
give the $j$th column of a matrix with $r\binom{d+m}{d}$ rows.  Divide
this column by $c_js^{\lambda\cdot e_j}$.  These normalized columns
still belong to $B^{\oplus r}$, and their reductions modulo $s$ are precisely the jets
of $Z^{e_j}u^{\gamma_j}$ at $(Q_i,U_i)$ in translated coordinates.
The assumed surjectivity supplies a maximal minor nonzero modulo $s$.
The original columns therefore span $B[1/s]^{\oplus r}$ over $F$.

We now interpret this rank statement as jet separation on a variety.
Setting $\delta=0$ in \eqref{eq:formal-substitution} gives actual
$F$-points $P_i$ of $Y^\circ\times\T^{d-h}$.  After inverting $s$, the
corresponding ring map factors through the order-$m$ jet algebra of
$(Y^\circ\times\T^{d-h})_F$ at $P_i$: elements nonvanishing at $P_i$
map to units, and its maximal
ideal maps into $(\delta)$.  Both this jet algebra and $B[1/s]$ have
dimension $\binom{d+m}{d}$ over $F$.  Surjectivity onto the direct sum
of the latter spaces forces each induced map from a jet algebra to be
an isomorphism.  It follows that the original functions separate the
order-$m$ jets at the $P_i$ themselves.  These points are distinct,
since two repeated centers would make two target blocks factor through
one jet algebra, contradicting surjectivity.

Thus the open jet-separation locus has an $F$-point and is nonempty.
The configuration space is irreducible, so this locus meets the
configuration space of any prescribed dense open subset.  A nonempty
open subset of a variety over the algebraically closed field $k$ has a
$k$-point.  This proves the assertion.  A local frame reduces the
line-bundle version to the one just proved.
\end{proof}

Together, Lemmas~\ref{lem:least-terms}--\ref{lem:formal-transfer} allow
us to replace finite-dimensional spaces by their least exponents, using
either of the two orders above, while preserving dimension and
transferring jet separation back to the original spaces.  Labels record
the exponents of variables that are left unchanged.

\subsection{Two elementary transformations}

\begin{lemma}\label{lem:unimodular}
For $A\in\operatorname{GL}_d(\Z)$, the transformation
$S_N\mapsto A S_N$ preserves cardinalities and gradedness, and has
backward jet transfer.
\end{lemma}

\begin{proof}
An invertible integral linear map preserves cardinalities and addition.
Its action on characters is induced by a torus automorphism, which
identifies the corresponding jet evaluations.
\end{proof}

The other transformation uses orders of vanishing to place each row in an
interval starting at zero.  The resulting exponents need not be consecutive.

\begin{lemma}[Coordinate compression]\label{lem:compression}
Let $(S_N)$ be a graded family in $\Z^d$, and select one coordinate,
written $z$.  For each degree $N$ and each label $\gamma\in\Z^{d-1}$,
form its row space
\[
 W_{N,\gamma}
 =\Span_k\{z^a:(a,\gamma)\in S_N\}.
\]
Replace this row by the set of orders at $z=1$ of nonzero elements of
$W_{N,\gamma}$, retaining the label $\gamma$.  The resulting family
preserves cardinalities and gradedness, and has backward jet transfer.
If the original row lies in a real interval of length $b$, every new
exponent lies in $[0,b]$.
\end{lemma}

\begin{proof}
Expand the row spaces in the parameter $z-1$.  Lemma~\ref{lem:least-terms}
proves the cardinality and gradedness assertions, because multiplication
adds both degrees and labels.  Lemma~\ref{lem:formal-transfer}, with
$Y=\mathbb G_m$, $y=1$, and weight $1$, proves backward jet transfer.

For a nonempty row let $a_-$ and $a_+$ be its smallest and largest
exponents.  Multiplication by $z^{-a_-}$ preserves orders at $1$ and
turns each member into a polynomial of degree at most $a_+-a_-\leq b$.
Its order of vanishing is at most its degree, proving the bound.
\end{proof}

Consequently, suppose a family is bounded by dilates of a polygon.
Replacing each slice of that polygon parallel to the compressed coordinate
by an interval of the same length starting at zero gives a new region
whose dilates contain the compressed family.  This applies equally to real
slice endpoints and to every dilation of the polygon.  In characteristic
$p$, a compressed row need not contain all intermediate integers: for
example, $\Span_k\{1,z^p\}$ has orders $0$ and $p$ at $1$.  Only the
interval bound of Lemma~\ref{lem:compression} will be used.
